% Сборка: pdflatex -interaction=nonstopmode -halt-on-error multi_index_algorithm.tex
\documentclass[11pt,a4paper]{article}
\usepackage[T2A]{fontenc}
\usepackage[utf8]{inputenc}
\usepackage[english,russian]{babel}
\usepackage[margin=22mm]{geometry}
\usepackage{amsmath,amssymb,bm,microtype}
\usepackage[hidelinks]{hyperref}
\setlength{\parindent}{0pt}
\setlength{\parskip}{4pt}
\setlength{\emergencystretch}{2em}
\newcommand{\R}{\mathbb R}
\newcommand{\norm}[1]{\lVert#1\rVert}
\newcommand{\code}[1]{\texttt{\detokenize{#1}}}
\begin{document}
\begin{center}
{\Large\bfseries Обучение мультииндекса в ADP}\\[3pt]
{\large Текущий алгоритм и солвер Grassmann}\\[5pt]
{\small Реализация на 24 октября 2026 года}
\end{center}

\section{Модель и начальное приближение}

По наблюдениям $(x_i,y_i)_{i=1}^n$, $x_i\in\R^d$, оценивается модель
\[
 y_i=\psi(B_*^\top x_i)+\varepsilon_i,
 \qquad \mathbb E[\varepsilon_i\mid x_i]=0,
 \qquad B_*^\top B_*=I_m,\quad m<d.
\]
Размерность $m$ задаётся заранее. Результат обучения --- подпространство
$\operatorname{col}(B_*)$: базисы $B_*$ и $B_*Q$ при ортогональной $Q$
описывают один индекс. Функция $\psi$ отдельно не обучается.
Основание метода --- равенство
$\nabla f(x)=B_*\nabla\psi(B_*^\top x)$: градиенты регрессии лежат
в искомом подпространстве. Для его идентификации градиенты должны
возбуждать все $m$ направлений.

Далее описан стандартный режим: нормированные моменты
(\code{estimator="new"}), локальная инициализация, ортогональный тензор
локализации и CPU-солвер \code{grassman}.
В формулах базис $B$ имеет размер $d\times m$; внутри солвера он передаётся
транспонированным, в виде $m\times d$.

Из выборки без возвращения выбираются $J$ центров $c_j=x_{i_j}$;
по умолчанию $J=n$. Для инициализации строятся изотропные веса
$w^{\rm lin}_{ji}=K(\norm{x_i-c_j}^2/h_{\rm lin}^2)$.
Масштаб $h_{\rm lin}$ выбирается по условию средней массы
$J^{-1}\sum_{j,i}w^{\rm lin}_{ji}\ge N_{\rm lin}$, где обычно
$N_{\rm lin}=2d$. В каждом центре решается локальная линейная регрессия
\[
 (\widehat a_j,\widehat b_j)
 =\arg\min_{a,b}\left\{
 \sum_iw^{\rm lin}_{ji}[y_i-a-b^\top(x_i-c_j)]^2
 +\lambda_{\rm init}\norm{b}^2\right\}.
\]
Свободный член не штрафуется; $\lambda_{\rm init}=10^{-8}$ по умолчанию.
Первые $m$ правых сингулярных векторов матрицы со строками
$\sqrt{N_j^{\rm lin}}\,\widehat b_j^\top$,
$N_j^{\rm lin}=\sum_iw^{\rm lin}_{ji}$, дают $B_0$.
При недостаточном численном ранге обучение завершается ошибкой.

\section{Адаптивные окрестности и локальные моменты}

На внешнем шаге $k$ фиксируются базис $B_k$, спектральные веса
$\ell_k=(\ell_{k,1},\ldots,\ell_{k,m})$, масштаб $h_k$
и коэффициент анизотропии $\alpha_k\in[0,1]$.
Для $\delta_{ji}=x_i-c_j$ и $z_{ji}=B_k^\top\delta_{ji}$ веса равны
\begin{equation}\label{eq:weights}
 w_{ji}=K(t_{ji}),\qquad
 t_{ji}=\frac{\alpha_k^2(\norm{\delta_{ji}}^2-\norm{z_{ji}}^2)
                  +\sum_{a=1}^m\ell_{k,a}z_{ji,a}^2}{h_k^2},
 \qquad K(t)=\max(1-t^2,0).
\end{equation}
Аргумент $K$ уже является квадратом нормированного расстояния:
в изотропном случае вес равен
$\max(1-\norm{\delta_{ji}}^4/h_k^4,0)$.
При уменьшении $\alpha_k$ окрестность расширяется в направлениях,
ортогональных текущему индексу. Начальные значения
$\ell_0=\bm 1$ и $\alpha_0=1$ дают изотропную геометрию;
спектр инициализации не подставляется вместо $\ell_0$.

Для каждого центра вычисляются масса и нормированные веса
\[
 N_j=\sum_iw_{ji}>0,\qquad a_{ji}=w_{ji}/N_j,\qquad
 \mu_j=\sum_i a_{ji}x_i,\qquad \bar y_j=\sum_i a_{ji}y_i.
\]
Затем генерируются $p=N_\phi>m$ независимых единичных направлений
$\phi_{js}=v_{js}/\norm{v_{js}}$, $v_{js}\sim N(0,I_d)$.
По умолчанию они заново генерируются на каждом внешнем шаге.
Матрица $\Phi_j\in\R^{p\times d}$ содержит строки $\phi_{js}^\top$.
Нормированные локальные моменты имеют вид
\begin{align}
 \Sigma_j&=\sum_i a_{ji}(x_i-\mu_j)(x_i-\mu_j)^\top,
 &c_j^{xy}&=\sum_i a_{ji}(x_i-\mu_j)(y_i-\bar y_j),\nonumber\\
 U_j&=\Phi_j\Sigma_j\in\R^{p\times d},
 &I_j&=\Phi_j c_j^{xy}\in\R^p.\label{eq:moments}
\end{align}
Код вычисляет эти произведения без явного формирования $\Sigma_j$
и дополнительно центрирует проекции для численной устойчивости.
При локальном приближении
$f(x)\approx a_j+(Yg_j)^\top x$ получается уравнение
$I_j\approx U_jYg_j$ с единым базисом $Y$ для всех центров.

\section{Задача и работа солвера Grassmann}

На одном вызове солвера $U_j,I_j,N_j$ фиксированы. Решается задача
\begin{equation}\label{eq:objective}
 \boxed{\quad
 \min_{Y^\top Y=I_m}F(Y),\qquad
 F(Y)=\sum_{j=1}^J N_j\min_{g_j\in\R^m}
                       \norm{I_j-U_jYg_j}^2.\quad}
\end{equation}
Это взвешенная ошибка локальных моментов в случайных направлениях.
На каждом пробном $Y$ локальные коэффициенты заново исключаются:
\[
 M_j=U_jY,\qquad g_j=M_j^\dagger I_j,
 \qquad r_j=I_j-M_jg_j.
\]
Знак $\dagger$ означает решение наименьших квадратов минимальной нормы.
Для хорошо обусловленных систем с $m\le2$ используется QR с повторной
ортогонализацией; в остальных случаях --- SVD и контроль ранга.
Поскольку $F(YQ)=F(Y)$, оптимизация ведётся по подпространствам
многообразия Грассмана $\mathrm{Gr}(m,d)$.

\paragraph{Градиент и малый шаг Гаусса--Ньютона.}
В области постоянного полного локального ранга горизонтальный градиент равен
\begin{equation}\label{eq:gradient}
 G=(I_d-YY^\top)\left[-2\sum_jN_j U_j^\top r_jg_j^\top\right].
\end{equation}
Из SVD градиента выбирается небольшой набор направлений
$V\in\R^{d\times q}$, $V^\top Y=0$, $V^\top V=I_q$.
По умолчанию сохраняется не менее $95\%$ его спектральной энергии,
с ограничением $q\le\min(m,d-m)$.
Касательный шаг ищется как $Z=VK$, $K\in\R^{q\times m}$;
поэтому линейная подзадача содержит только $qm$ неизвестных.

Обозначим через $\mathcal J$ якобиан составного остатка
$r=(\sqrt{N_j}\,r_j)_j$ по координатам $\operatorname{vec}K$.
Он учитывает зависимость $g_j$ от $Y$.
При полном ранге для приращения $D_j=U_jVK$ выполняется
\[
 \mathrm dg_j=(M_j^\top M_j)^{-1}
       (D_j^\top r_j-M_j^\top D_jg_j),\qquad
 \mathrm dr_j=-D_jg_j-M_j\,\mathrm dg_j.
\]
Обратная матрица здесь --- математическая запись: в коде применяются
QR/SVD-факторы. Демпфированная подзадача решается через
дополненную систему наименьших квадратов:
\begin{equation}\label{eq:gn}
 K=\arg\min_K\left\{
   \norm{\mathcal J\operatorname{vec}K+r}^2
   +\lambda\norm{K}_F^2\right\}.
\end{equation}
Обычно $\lambda=\code{lambda_penalty}=0.05$.
Это демпфирование численного шага. Отдельный параметр солвера
\code{local_ridge} добавляет штраф на $g_j$ в~\eqref{eq:objective};
по умолчанию он равен нулю и отличается от ridge инициализации.

\paragraph{Принятие шага и вырожденные случаи.}
Ортонормальность сохраняет полярная ретракция
\[
 Y_+=(Y+VK)(I_m+K^\top K)^{-1/2}.
\]
Длина $K$ ограничивается условием $\norm{K}_2\le\tan(0.5)$.
При поиске шага $K$ последовательно уменьшается вдвое,
а полный профиль $F(Y_+)$ пересчитывается.
Шаг принимается только при достаточном уменьшении по Армихо;
при неудаче демпфирование увеличивается.

Около потери локального ранга режим \code{relative_gn=True} может строить
предложение по регулярным центрам, если у остальных
$\operatorname{rank}(U_jY)=\operatorname{rank}(U_j)<m$.
Тогда демпфирование масштабируется как
$\lambda\norm{\mathcal J}_F^2/(qm)$.
Приём шага всё равно проверяется по полному функционалу.
Если такое предложение недоступно, используется вращение ранга один
с повторным решением локальных задач, до трёх квадратичных проб
и уменьшением шага. При превышении лимита рабочей памяти якобиана
(обычно $16$ МиБ) применяется спектральное направление спуска.
Центры с $U_j=0$ исключаются из поиска; их постоянный вклад
возвращается в значение ошибки.

На вызов солвера отводится до $20$ итераций. Проверяется условие
\[
 \norm{G}_F\le10^{-6}\max\left(1,\sum_jN_j\norm{I_j}^2\right).
\]
Достижение бюджета или неудача поиска шага также останавливают солвер.
На границе ранга малый градиент не считается сертификатом гладкой
стационарности. Возвращаются базис, коэффициенты и диагностика:
ошибка, норма градиента, потеря ранга и причина остановки.

\section{Внешнее обновление и итог обучения}

После солвера формируется малая матрица
$C=\sum_jN_jg_jg_j^\top$.
Если $C=Q\operatorname{diag}(\gamma_1,\ldots,\gamma_m)Q^\top$,
$\gamma_1\ge\cdots\ge\gamma_m>0$, то
\[
 B_{k+1}=YQ,\qquad \ell_{k+1,a}=\gamma_a/\gamma_1.
\]
В реализации это SVD матрицы со строками $\sqrt{N_j}\,g_j^\top$.
Поворот сохраняет подпространство; коэффициенты поворачиваются
согласованно. Недостаточный ранг коэффициентов приводит к ошибке.

Начальный $h_0\ge h_{\min}$ выбирается бисекцией по условию
$J^{-1}\sum_jN_j\ge N_{\rm loc}$ в изотропной геометрии.
После каждого шага предлагается $h_{k+1}=h_k/a$, обычно $a=\sqrt2$.
Для нового базиса и спектра бисекцией выбирается наибольшее
$\alpha_{k+1}\in[0,1]$, сохраняющее среднюю массу не меньше
$N_{\rm loc}$. Это условие на сумму весов, а не на число соседей.
Внешний цикл заканчивается, когда следующий $h$ ниже $h_{\min}$,
масса недостижима даже при $\alpha=0$ либо исчерпан \code{outer_steps}.
Внутренняя сходимость солвера сама по себе внешний цикл не завершает.

По умолчанию \code{select_step="best"} возвращает базис внешнего шага
с наименьшей ошибкой
\[
 E_k=\sum_{j=1}^J(y_{i_j}-\bar y_{j,k})^2.
\]
Здесь $\bar y_{j,k}$ --- локальное среднее из статистик,
построенных по \emph{входящей} геометрии шага $k$.
Оценка $E_k$ относится к соответствующему выходному базису $B_{k+1}$
и отличается от $F(Y)$.
При стандартном \code{training_set="all"} эти же наблюдения участвуют
в обучении: независимой валидации здесь нет.
Режим \code{select_step="last"} возвращает последний базис.

Внешние шаги меняют окрестности и случайные направления, поэтому
единого фиксированного функционала у всего цикла нет.
Убывание внутренней ошибки не гарантирует улучшения восстановления
истинного подпространства или достижения глобального минимума.
Итоговая матрица доступна как \code{model.basis_};
\code{model.transform(X)} вычисляет $XB$.

\paragraph{Параметры и исходники.}
Другие стандартные значения: $N_{\rm loc}=10$,
$p=\min(N_{\rm loc},d)$ с обязательным $p>m$,
$h_{\min}=1$, \code{seed=42}.
Описание следует файлам
\path{ADP/core/multi/ADP_multi_index.py},
\path{ADP/engine/common/index_fit.py},
\path{ADP/engine/common/initialize.py},
\path{ADP/engine/common/weights.py},
\path{ADP/engine/common/statistic.py}
и \path{ADP/solver/grassman.py}.

\end{document}
